% !TEX program = pdflatex
\documentclass[11pt,a4paper]{article}

\usepackage[margin=1in]{geometry}
\usepackage{parskip}

\usepackage{amsmath,amssymb,amsthm}
\usepackage{xcolor}

\theoremstyle{definition}
\newtheorem{definition}{Definition}
\newtheorem{example}{Example}

\theoremstyle{plain}
\newtheorem{theorem}{Theorem}
\newtheorem{proposition}{Proposition}
\newtheorem{lemma}{Lemma}

\usepackage{mathtools,longtable,booktabs,array}
\setlength{\emergencystretch}{2em}
\allowdisplaybreaks[1]
\raggedbottom
\begin{document}


\begin{center}
    {\LARGE \textbf{Mathematical Methods for Physics--Lecture 15}} \\[0.6em]
    {\large Convolution, Fourier Applications, and Laplace Transforms}
\end{center}

\vspace{-0.5em}
\noindent\rule{\textwidth}{0.5pt}
\vspace{0.3em}

\begin{center}
    Instructor: Zhaolong Han, SOCE, Shanghai Jiao Tong University;
\end{center}

\vspace{-0.6em}
\noindent Email: han.arkey@sjtu.edu.cn

\vspace{0.4em}
\noindent\rule{\textwidth}{0.5pt}

\vspace{0.8em}

\section{Convolution and Fourier Transforms}

\subsection{Definition and Algebraic Properties of Convolution}

\begin{definition}[Convolution]

For two functions \(f_1(t)\) and \(f_2(t)\) defined on
\((-\infty,+\infty)\), their convolution is defined by
\[
(f_1*f_2)(t)
=\int_{-\infty}^{+\infty}
f_1(\tau)f_2(t-\tau)\,\mathrm d\tau.
\]

\end{definition}

\textbf{Algebraic properties}
\[
f_1*f_2=f_2*f_1,
\]
\[
f_1*(f_2*f_3)=(f_1*f_2)*f_3,
\]
\[
f_1*(f_2+f_3)=f_1*f_2+f_1*f_3.
\]

\subsection{The Convolution Theorem}

Let
\[
\mathcal F[f_1(t)]=F_1(\omega),
\qquad
\mathcal F[f_2(t)]=F_2(\omega).
\]

\textbf{(1) The Fourier transform of a convolution is the product of the
transforms}
\[
\mathcal F[(f_1*f_2)(t)]
=\mathcal F[f_1]\cdot\mathcal F[f_2]
=F_1(\omega)F_2(\omega).
\]

That is,
\[
f_1*f_2
\ \xleftrightarrow{\mathcal F}\
F_1F_2.
\]

\textbf{(2) Fourier transform of a product}
\[
\mathcal F[f_1f_2]
=\frac1{2\pi}\bigl(\mathcal F[f_1]*\mathcal F[f_2]\bigr)
=\frac1{2\pi}(F_1*F_2).
\]

\textbf{Proof of (1)}\par\nopagebreak[4]
\[
\begin{aligned}
\mathcal F[(f_1*f_2)(t)]
&=\int_{-\infty}^{+\infty}
(f_1*f_2)(t)e^{-i\omega t}\,\mathrm dt\\
&=\int_{-\infty}^{+\infty}
\left[
\int_{-\infty}^{+\infty}
f_1(\tau)f_2(t-\tau)\,\mathrm d\tau
\right]e^{-i\omega t}\,\mathrm dt\\
&=\int_{-\infty}^{+\infty}
f_1(\tau)
\left[
\int_{-\infty}^{+\infty}
f_2(t-\tau)e^{-i\omega t}\,\mathrm dt
\right]\mathrm d\tau.
\end{aligned}
\]

By the shifting property,
\[
\int_{-\infty}^{+\infty}
f_2(t-\tau)e^{-i\omega t}\,\mathrm dt
=e^{-i\tau\omega}F_2(\omega).
\]

Therefore,
\[
\begin{aligned}
\mathcal F[(f_1*f_2)(t)]
&=\int_{-\infty}^{+\infty}
f_1(\tau)e^{-i\tau\omega}F_2(\omega)\,\mathrm d\tau\\
&=F_2(\omega)
\int_{-\infty}^{+\infty}
f_1(\tau)e^{-i\tau\omega}\,\mathrm d\tau\\
&=F_2(\omega)F_1(\omega).
\end{aligned}
\]

\subsection{Example: Proving the Integration Property by Convolution}

If
\[
\int_{-\infty}^{+\infty}f(t)\,\mathrm dt\ne0,
\]

then
\[
\mathcal F\!\left[
\int_{-\infty}^{t}f(\tau)\,\mathrm d\tau
\right]
=\frac{F(\omega)}{i\omega}
+\pi\delta(\omega)F(0).
\]

\textbf{Proof}\par\nopagebreak[4]
First, we show that
\[
\int_{-\infty}^{t}f(\tau)\,\mathrm d\tau
=(f*u)(t),
\]

where
\[
u(t)=
\begin{cases}
1,&t>0,\\
0,&t<0.
\end{cases}
\]

Indeed,
\[
(f*u)(t)
=\int_{-\infty}^{+\infty}
f(\tau)u(t-\tau)\,\mathrm d\tau.
\]

Let \(t-\tau=\xi\). Then
\[
\begin{aligned}
(f*u)(t)
&=\int_{+\infty}^{-\infty}
f(t-\xi)u(\xi)(-\mathrm d\xi)\\
&=\int_{-\infty}^{+\infty}
f(t-\xi)u(\xi)\,\mathrm d\xi\\
&=\int_{-\infty}^{0}0\,\mathrm d\xi
+\int_0^{+\infty}f(t-\xi)\,\mathrm d\xi.
\end{aligned}
\]

Now let \(t-\xi=\tau\). We obtain
\[
(f*u)(t)
=\int_t^{-\infty}f(\tau)(-\mathrm d\tau)
=\int_{-\infty}^{t}f(\tau)\,\mathrm d\tau.
\]

By the convolution theorem,
\[
\begin{aligned}
\mathcal F\!\left[
\int_{-\infty}^{t}f(\tau)\,\mathrm d\tau
\right]
&=\mathcal F[(f*u)(t)]\\
&=\mathcal F[f]\cdot\mathcal F[u]\\
&=F(\omega)
\left[\frac1{i\omega}+\pi\delta(\omega)\right]\\
&=\frac{F(\omega)}{i\omega}
+F(\omega)\pi\delta(\omega)\\
&=\frac{F(\omega)}{i\omega}
+F(0)\pi\delta(\omega).
\end{aligned}
\]

\section{Applications of Fourier Transforms}

Fourier transforms can be used for

\begin{enumerate}
\item
  Solving differential equations.

\item
  Solving integral equations.

\end{enumerate}

\subsection{Example 7.20: Solving a Differential Equation}

Solve
\[
x''(t)-x(t)=-f(t),
\qquad -\infty<t<+\infty.
\]

Assume the boundary conditions
\[
x(\pm\infty)=0,
\qquad
x'(\pm\infty)=0.
\]

\textbf{Solution}\par\nopagebreak[4]
\textbf{Step 1. Take the Fourier transform.}\par\nopagebreak[4]
Let
\[
\mathcal F[x(t)]=\overline X(\omega),
\qquad
\mathcal F[f(t)]=F(\omega).
\]

Taking the Fourier transform of both sides gives
\[
\mathcal F[x''(t)-x(t)]=-\mathcal F[f(t)].
\]

By the differentiation property,
\[
(i\omega)^2\overline X(\omega)-\overline X(\omega)
=-F(\omega),
\]

Therefore,
\[
\overline X(\omega)=\frac{F(\omega)}{\omega^2+1}.
\]

\textbf{Step 2. Take the inverse Fourier transform.}\par\nopagebreak[4]
\[
\begin{aligned}
x(t)
&=\mathcal F^{-1}[\overline X(\omega)]\\
&=\mathcal F^{-1}\!\left[
F(\omega)\frac1{\omega^2+1}
\right]\\
&=\mathcal F^{-1}[F(\omega)]
*\mathcal F^{-1}\!\left[\frac1{\omega^2+1}\right]\\
&=f(t)*\left(\frac12e^{-|t|}\right)\\
&=\frac12\int_{-\infty}^{+\infty}
f(\tau)e^{-|t-\tau|}\,\mathrm d\tau.
\end{aligned}
\]

The transform pair used here is
\[
e^{-\beta|t|}
\ \xleftrightarrow{\mathcal F}\
\frac{2\beta}{\omega^2+\beta^2},
\]

It follows that
\[
\mathcal F^{-1}\!\left[\frac1{\omega^2+\beta^2}\right]
=\frac1{2\beta}e^{-\beta|t|},
\qquad \beta>0.
\]

\subsection{Example 7.22: Solving a Convolution Integral Equation}

Solve the integral equation
\[
\int_{-\infty}^{+\infty}
e^{-|t-\xi|}x(\xi)\,\mathrm d\xi
=f(t),
\]

where \(f(t)\) is known.

\textbf{Analysis.} The left-hand side is a convolution, so the equation
can be written as
\[
e^{-|t|}*x(t)=f(t).
\]

\textbf{Step 1. Take the Fourier transform.}\par\nopagebreak[4]
By the convolution theorem,
\[
\mathcal F[e^{-|t|}*x(t)]
=\mathcal F[f(t)],
\]
\[
\mathcal F[e^{-|t|}]\cdot\mathcal F[x(t)]
=F(\omega).
\]

That is,
\[
\frac2{\omega^2+1}\overline X(\omega)=F(\omega).
\]

Therefore,
\[
\overline X(\omega)
=\frac{\omega^2+1}{2}F(\omega).
\]

\textbf{Step 2. Take the inverse Fourier transform.}\par\nopagebreak[4]
\[
\begin{aligned}
x(t)
&=\mathcal F^{-1}\!\left[
\frac{\omega^2+1}{2}F(\omega)
\right]\\
&=\frac12\left\{
\mathcal F^{-1}[\omega^2F(\omega)]
+\mathcal F^{-1}[F(\omega)]
\right\}\\
&=\frac12\left\{
-\mathcal F^{-1}[(i\omega)^2F(\omega)]
+f(t)
\right\}\\
&=\frac12[-f''(t)+f(t)].
\end{aligned}
\]

The differentiation property used here is
\[
f''(t)\ \xleftrightarrow{\mathcal F}\
(i\omega)^2F(\omega),
\]

with the boundary conditions
\[
f(\pm\infty)=0,
\qquad
f'(\pm\infty)=0.
\]

\textbf{In particular}, let
\[
f(t)=e^{-|t|}-\frac12e^{-2|t|},
\]

Substituting into the formula above gives
\[
x(t)=\frac34e^{-2|t|}.
\]

\subsection{Example 7.21: Solving an Integral Equation with a Sine Kernel}

Solve
\[
\int_0^{+\infty}x(\omega)\sin\omega t\,\mathrm d\omega
=e^{-t},
\qquad t>0.
\]

\subsubsection*{Review: The Fourier Integral Formula (Inverse Transform)}
\[
f(t)=\frac1{2\pi}
\int_{-\infty}^{+\infty}
\left[
\int_{-\infty}^{+\infty}
f(\tau)e^{-i\omega\tau}\,\mathrm d\tau
\right]e^{i\omega t}\,\mathrm d\omega.
\]

\textbf{Special cases}

\textbf{(I) Even functions}

If \(f(-t)=f(t)\), then
\[
f(t)=\frac2\pi
\int_0^{+\infty}
\left[
\int_0^{+\infty}
f(\tau)\cos\omega\tau\,\mathrm d\tau
\right]\cos\omega t\,\mathrm d\omega.
\]

This is called the \textbf{Fourier cosine integral formula}.

\textbf{(II) Odd functions}

If \(f(-t)=-f(t)\), then
\[
f(t)=\frac2\pi
\int_0^{+\infty}
\left[
\int_0^{+\infty}
f(\tau)\sin\omega\tau\,\mathrm d\tau
\right]\sin\omega t\,\mathrm d\omega.
\]

This is called the \textbf{Fourier sine integral formula}.

\subsubsection*{Return to the Integral Equation}

From
\[
\int_0^{+\infty}x(\omega)\sin\omega t\,\mathrm d\omega=e^{-t},
\]

we obtain
\[
\frac2\pi\int_0^{+\infty}
x(\omega)\sin\omega t\,\mathrm d\omega
=\frac2\pi e^{-t}.
\]

The Fourier sine integral formula gives
\[
x(\omega)=\int_0^{+\infty}
\frac2\pi e^{-t}\sin\omega t\,\mathrm dt.
\]

The integral can be evaluated in either of two ways.

\begin{enumerate}
\item
  Method 1: integrate by parts twice.

\item
  Method 2: use Euler's formula,

\end{enumerate}
\[
\sin\omega t=\frac{e^{i\omega t}-e^{-i\omega t}}{2i}.
\]

Both methods give
\[
\int_0^{+\infty}e^{-t}\sin\omega t\,\mathrm dt
=\frac{\omega}{1+\omega^2},
\qquad
x(\omega)=\frac{2\omega}{\pi(1+\omega^2)}.
\]

\section{The Laplace Transform}

\subsection{Definition of the Laplace Transform}

\begin{definition}[Laplace transform]
For a function \(f(t)\) defined on \([0,+\infty)\), define
\[
F(p)=\int_0^{+\infty}f(t)e^{-pt}\,\mathrm dt,
\]

where
\[
p=\sigma+i\omega
\]

is a complex variable.

The pair \(\{f(t),F(p)\}\) is called a Laplace transform pair. We write
\[
\mathcal L[f(t)]=F(p),
\]

or
\[
\mathcal L^{-1}[F(p)]=f(t).
\]

\end{definition}

\subsection{Functions of Exponential Order and the Existence Theorem}

\begin{definition}[Exponential order]
If there exist constants \(M>0\) and \(\sigma_c\ge0\) such that
\[
|f(t)|\le Me^{\sigma_ct},
\]

then \(f(t)\) is said to be of exponential order, and \(\sigma_c\) is
its growth exponent. The bound is assumed for \(t\ge0\).
\end{definition}

\begin{theorem}[Existence of the Laplace transform]

Suppose that \(f(t)\), defined on \([0,+\infty)\), satisfies the
following conditions.

\begin{enumerate}
\def\labelenumi{\arabic{enumi}.}
\item
  \(f(t)\) satisfies the Dirichlet conditions on every finite interval.

\item
  \(f(t)\) is of exponential order, that is,
  \(|f(t)|\le Me^{\sigma_ct}\).

\end{enumerate}

Then the Laplace transform
\[
F(p)=\int_0^{+\infty}f(t)e^{-pt}\,\mathrm dt
\]

exists and is analytic in the right half-plane
\[
\operatorname{Re}p>\sigma_c.
\]
\end{theorem}

\subsubsection*{Proof: Absolute Convergence}

Since
\[
|f(t)|\le Me^{\sigma_ct}
\]

and
\[
|e^{-pt}|
=|e^{-\sigma t}e^{-i\omega t}|
=e^{-\sigma t},
\]

for \(\sigma>\sigma_c\) we have
\[
\begin{aligned}
\int_0^{+\infty}|f(t)e^{-pt}|\,\mathrm dt
&\le\int_0^{+\infty}Me^{\sigma_ct}|e^{-pt}|\,\mathrm dt\\
&=\int_0^{+\infty}Me^{-(\sigma-\sigma_c)t}\,\mathrm dt\\
&=\left.
M\frac{e^{-(\sigma-\sigma_c)t}}{-(\sigma-\sigma_c)}
\right|_0^{+\infty}\\
&=\frac{M}{\sigma-\sigma_c}.
\end{aligned}
\]

Therefore, the improper integral
\[
\int_0^{+\infty}f(t)e^{-pt}\,\mathrm dt
\]

exists and converges absolutely.

\subsubsection*{Proof: Analyticity}

Both the defining integral and
\[
\int_0^{+\infty}
f(t)\frac{\partial}{\partial p}(e^{-pt})\,\mathrm dt
\]

converge locally uniformly in \(\operatorname{Re}p>\sigma_c\), allowing
differentiation with respect to the parameter \(p\).
\[
\begin{aligned}
F'(p)
&=\frac{\mathrm d}{\mathrm dp}
\int_0^{+\infty}f(t)e^{-pt}\,\mathrm dt\\
&=\int_0^{+\infty}
f(t)\frac{\partial}{\partial p}(e^{-pt})\,\mathrm dt\\
&=\int_0^{+\infty}(-t)f(t)e^{-pt}\,\mathrm dt.
\end{aligned}
\]

Hence \(F(p)\) is analytic in \(\operatorname{Re}p>\sigma_c\).

\subsection{Common Laplace Transform Pairs}

\subsubsection*{The Unit Step Function}
\[
u(t)\ \xleftrightarrow{\mathcal L}\ \frac1p,
\qquad \operatorname{Re}p>0.
\]

Since
\[
\begin{aligned}
\mathcal L[u(t)]
&=\int_0^{+\infty}1\cdot e^{-pt}\,\mathrm dt\\
&=\left.\frac{e^{-pt}}{-p}\right|_0^{+\infty}\\
&=\frac1p.
\end{aligned}
\]

\subsubsection*{Power Functions}

Similarly,
\[
\begin{aligned}
\mathcal L[t]
&=\int_0^{+\infty}te^{-pt}\,\mathrm dt\\
&=\int_0^{+\infty}t\,\mathrm d\!\left(\frac{e^{-pt}}{-p}\right)\\
&=\left.-\frac{te^{-pt}}p\right|_0^{+\infty}
+\int_0^{+\infty}\frac{e^{-pt}}p\,\mathrm dt\\
&=0+\left.\frac{e^{-pt}}{-p^2}\right|_0^{+\infty}\\
&=\frac1{p^2}.
\end{aligned}
\]
\[
\begin{aligned}
\mathcal L[t^2]
&=\int_0^{+\infty}t^2e^{-pt}\,\mathrm dt\\
&=\left.t^2\frac{e^{-pt}}{-p}\right|_0^{+\infty}
-\int_0^{+\infty}\frac{e^{-pt}}{-p}\,2t\,\mathrm dt\\
&=\frac2p\cdot\frac1{p^2}\\
&=\frac{2!}{p^3}.
\end{aligned}
\]

Generalizing gives
\[
\mathcal L[t^n]=\frac{n!}{p^{n+1}},
\]

That is,
\[
t^n\ \xleftrightarrow{\mathcal L}\
\frac{n!}{p^{n+1}},
\qquad n=0,1,2,\ldots,\quad \operatorname{Re}p>0.
\]

\subsubsection*{An Exponential Function}
\[
e^{at}\ \xleftrightarrow{\mathcal L}\ \frac1{p-a},
\]

where \(a\) is a constant.
\[
\begin{aligned}
\mathcal L[e^{at}]
&=\int_0^{+\infty}e^{at}e^{-pt}\,\mathrm dt\\
&=\left.\frac{e^{(a-p)t}}{a-p}\right|_0^{+\infty}\\
&=\frac1{p-a},
\end{aligned}
\]

The convergence condition is
\[
\operatorname{Re}(a-p)<0
\quad\Longleftrightarrow\quad
\operatorname{Re}p>\operatorname{Re}a.
\]

\subsubsection*{Trigonometric and Hyperbolic Functions}

For real \(\beta\), the trigonometric pairs converge for \(\operatorname{Re}p>0\);
the hyperbolic pairs converge for \(\operatorname{Re}p>|\beta|\).

\begingroup
\small
\renewcommand{\arraystretch}{1.65}
\everymath{\displaystyle}
\begin{longtable}{@{}p{0.29\textwidth}p{0.65\textwidth}@{}}
\toprule
Function & Laplace transform \\
\midrule
\endhead
\bottomrule
\endlastfoot
\(\sin\beta t\) & \(\dfrac{\beta}{p^2+\beta^2}\) \\
\(\cos\beta t\) & \(\dfrac{p}{p^2+\beta^2}\) \\
\(\sinh\beta t\) & \(\dfrac{\beta}{p^2-\beta^2}\) \\
\(\cosh\beta t\) & \(\dfrac{p}{p^2-\beta^2}\) \\
\end{longtable}
\endgroup

\textbf{Derivation for the sine function}
\[
\begin{aligned}
\mathcal L[\sin\beta t]
&=\mathcal L\!\left[
\frac{e^{i\beta t}-e^{-i\beta t}}{2i}
\right]\\
&=\frac1{2i}\left\{
\mathcal L[e^{i\beta t}]
-\mathcal L[e^{-i\beta t}]
\right\}\\
&=\frac1{2i}\left[
\frac1{p-i\beta}-\frac1{p+i\beta}
\right]\\
&=\frac1{2i}\cdot\frac{2i\beta}{p^2+\beta^2}\\
&=\frac{\beta}{p^2+\beta^2}.
\end{aligned}
\]

\section{Properties of the Laplace Transform}

\subsection{Linearity}
\[
\mathcal L[af_1+bf_2]
=a\mathcal L[f_1]+b\mathcal L[f_2].
\]
\[
\mathcal L^{-1}[aF_1+bF_2]
=a\mathcal L^{-1}[F_1]+b\mathcal L^{-1}[F_2].
\]

\subsubsection*{Example: Find the Inverse Laplace Transform}
\[
\mathcal L^{-1}\!\left[
\frac{p}{(p-1)(p^2+4)}
\right].
\]

First, decompose the expression into partial fractions.
\[
\begin{aligned}
\frac{p}{(p-1)(p^2+4)}
&=\frac{A}{p-1}+\frac{Bp+C}{p^2+4}\\
&=\frac{1/5}{p-1}
+\frac{-p/5+4/5}{p^2+4}.
\end{aligned}
\]

Therefore,
\[
\begin{aligned}
\mathcal L^{-1}\!\left[
\frac{p}{(p-1)(p^2+4)}
\right]
&=\frac15\mathcal L^{-1}\!\left[\frac1{p-1}\right]
-\frac15\mathcal L^{-1}\!\left[\frac{p}{p^2+4}\right]\\
&\quad+\frac45\mathcal L^{-1}\!\left[\frac1{p^2+4}\right]\\
&=\frac15e^t-\frac15\cos2t+\frac25\sin2t.
\end{aligned}
\]

The transform pairs used are
\[
e^{at}\ \xleftrightarrow{\mathcal L}\ \frac1{p-a},
\]
\[
\cos at\ \xleftrightarrow{\mathcal L}\ \frac{p}{p^2+a^2},
\]
\[
\sin at\ \xleftrightarrow{\mathcal L}\ \frac{a}{p^2+a^2}.
\]

\subsection{Scaling}

Let
\[
f(t)\ \xleftrightarrow{\mathcal L}\ F(p).
\]

then
\[
f(at)\ \xleftrightarrow{\mathcal L}\
\frac1aF\!\left(\frac pa\right),
\qquad a>0.
\]

and
\[
\frac1a f\!\left(\frac ta\right)
\ \xleftrightarrow{\mathcal L}\
F(ap).
\]

\subsection{Time Delay}

For \(t_0\ge0\),
\[
f(t-t_0)u(t-t_0)
\ \xleftrightarrow{\mathcal L}\
e^{-pt_0}F(p).
\]

That is,
\[
\mathcal L[f(t-t_0)u(t-t_0)]
=e^{-pt_0}F(p),
\]

and
\[
\mathcal L^{-1}[e^{-pt_0}F(p)]
=f(t-t_0)u(t-t_0).
\]

\subsection{Shifting in the Transform Variable}
\[
e^{p_0t}f(t)
\ \xleftrightarrow{\mathcal L}\
F(p-p_0).
\]

\textbf{Proof.} From the definition,
\[
\begin{aligned}
\mathcal L[e^{p_0t}f(t)]
&=\int_0^{+\infty}f(t)e^{p_0t}e^{-pt}\,\mathrm dt\\
&=\int_0^{+\infty}f(t)e^{-(p-p_0)t}\,\mathrm dt\\
&=F(p-p_0).
\end{aligned}
\]
\end{document}
